\documentclass[../../infufrgs/ppgc-prop-tese.tex]{subfiles}
\begin{document}
% \begin{table}[ht]

\small\begin{longtable}{p{0.13\textwidth}p{0.38\textwidth}p{0.38\textwidth}}
\centering
\textbf{Month} & \textbf{Thesis writing and structure} & \textbf{Experiments and analysis} \\
\hline
Oct 2025 &
\begin{itemize}
  \item Define a detailed thesis outline (chapters and sections).
  \item Map how existing papers (grokking, logical steering, LSNC) become thesis chapters.
  \item Draft the initial version of Chapter~1 (Introduction): motivation, problem statement, overview of contributions.
\end{itemize}
&
\begin{itemize}
  \item Organize code, data and experiment folders for grokking and steering.
  \item Decide which experimental results are final and which need to be rerun or cleaned up.
  \item Align naming and file structure for figures and tables to be reused in the thesis.
\end{itemize}
\\

Nov 2025 &
\begin{itemize}
  \item Define a detailed thesis outline (chapters and sections).
  \item Map how existing papers (grokking, logical steering, LSNC) become thesis chapters.
  \item Draft the initial version of Chapter~1 (Introduction): motivation, problem statement, overview of contributions.
\end{itemize}
&
\begin{itemize}
  \item Organize code, data and experiment folders for grokking and steering.
  \item Decide which experimental results are final and which need to be rerun or cleaned up.
  \item Align naming and file structure for figures and tables to be reused in the thesis.
\end{itemize}
\\

Dec 2025 &
\begin{itemize}
  \item Turn the grokking paper into its thesis chapter: full narrative, methodology and results.
  \item Write discussion and threats-to-validity subsections for the grokking chapter.
  \item Draft the structure of Chapter~2 (Background): transformers, interpretability, grokking, soft logic (LTN/PSL), etc.
\end{itemize}
&
\begin{itemize}
  \item Make the grokking repository “thesis-ready” (README, tags, clear experiment scripts).
  \item Decide what will be kept from LSNC as a “negative result / lessons learned” section.
  \item Collect and centralize logs and results related to LSNC that will appear in the thesis.
\end{itemize}
\\

Jan 2026 &
\begin{itemize}
  \item Define the structure of the steering chapter (method, LSNC trajectory / negative result, final approach).
  \item Write the method subsection for logical validators and reranking (including mathematical formulation and intuition).
  \item Draft the results section for steering, including placeholders for single-rule and multi-rule scenarios.
\end{itemize}
&
\begin{itemize}
  \item Finalize the steering pipeline: validators, logit mapping, integration with decoding.
  \item Run main experiments for single-rule steering and compare with baselines (FUDGE, GeDi, PPLM, DExperts, etc.).
  \item Run multi-rule experiments and ablations (with/without certain rules, with/without calibration, etc.).
  \item Generate final plots and tables for the steering chapter in a format ready for the thesis.
\end{itemize}
\\

Feb 2026 &
\begin{itemize}
  \item Integrate and revise the steering chapter so that notation and references match the background and grokking chapters.
  \item General discussion / synthesis chapter (e.g., Chapter~5): how grokking and steering fit together, limitations, future work.
  \item Review of Chapter~1 (Introduction) and the conclusions so that they reflect the final set of contributions.
  \item First complete version of the thesis.
\end{itemize}
&
\begin{itemize}
  \item Check that the full thesis compiles cleanly (no missing references or figures).
  \item Minimal reproducibility material (scripts or notebooks) for the main results.
  \item Defense slides: story arc, key figures, and backup slides for extra results.
\end{itemize}
\\

Mar 2026 &
\begin{itemize}
  \item Incorporate advisor feedback into all chapters.
  \item Polish the text: transitions, reduction of repetition, consistent terminology and notation.
  \item Finalize pre-textual elements: abstract, summary, acknowledgements, lists of figures/tables, etc.
  \item Freeze the final defense version of the thesis.
\end{itemize}
&

\\
\hline
\caption{Monthly plan for completing the PhD thesis (Oct 2025 -- Mar 2026).}
\end{longtable}
% \end{table}





\begin{table}[ht]
\centering
\small
\begin{tabular}{p{0.32\textwidth}p{0.58\textwidth}}
\hline
\multicolumn{2}{l}{\textbf{Completed items}} \\
%\hline
\textbf{Type} & \textbf{Description} \\

\hline
Research direction and scope defined &
Overall PhD topic, main research questions, and methodological approach are clearly defined (grokking + logic-based steering of neural models). \\

Conference paper accepted (requirement (ii)) &
At least one full paper accepted in a venue classified as A3 or higher (e.g., ICTAI conference paper \emph{“Inducing Grokking with Distribution Shifts”}, subject to confirmation of Qualis level by the program). \\

Core grokking experiments and analysis &
Main grokking experiments are designed, implemented and validated; results are consolidated in the conference paper and ready to be adapted into a thesis chapter. \\

Preliminary logic-based steering work &
Initial implementations and experiments with logic-based steering (e.g., LSNC / soft-logic approaches) performed, providing material for a “trajectory / lessons learned” part of the thesis. \\
\hline
\vspace{1cm}
\end{tabular}

\begin{tabular}{p{0.32\textwidth}p{0.58\textwidth}}
\multicolumn{2}{l}{\textbf{Items to be completed}} \\
\textbf{Type} & \textbf{Description} \\
\hline
Journal submission (requirement (i)) &
Submition of a  full paper to a journal classified as Qualis A3 or higher, derived from the thesis results. \\

Final steering method and experiments &
Consolidate the final version of the logic-based steering / reranking method, run the definitive experiments, and prepare final tables/figures. \\

Thesis chapters – writing &
Complete writing of all thesis chapters (Introduction, Background, Grokking chapter, Steering chapter, Discussion/Conclusions), integrating and polishing existing material from the papers. \\

Thesis compilation and internal review &
Produce a full, coherent thesis manuscript; ensure it compiles correctly, check cross-references and formatting, and obtain advisor’s feedback on the complete draft. \\

Administrative and formal requirements &
Meet all formal program requirements: deliver printed/digital copies, fill in forms, meet deadlines, follow the official template, and schedule the defense according to regulations. \\

Defense preparation and final corrections &
Prepare and rehearse the defense presentation, defend the thesis, and implement any mandatory post-defense corrections before the final deposit. \\
\hline
\end{tabular}
\caption{Summary of completed and to be completed PhD items.}
\end{table}


\end{document}